\section{Conclusion}
\label{sec:conclusion}

We presented AnchorLoRA, a modality-asymmetric adaptation strategy for cross-domain multimodal remote sensing class-incremental learning. By leveraging the empirical observation that spectral features are inherently more stable across geographic domains than spatial features, AnchorLoRA freezes the spectral branch as a semantic anchor and adapts spatial branches with low-rank adapters whose capacity matches each modality's measured drift. On the Marathon CIL benchmark spanning 3 domains, 9 tasks, and 32 classes, AnchorLoRA achieves 75.7\% average accuracy---the highest among all exemplar-free methods tested.

\paragraph{Limitations.}
The warm-up phase length (3 tasks) coincides with the first domain's task count, though ablation shows warmup=2 performs comparably ($-$1.1pp). Houston domain accuracy degrades substantially ($-$24.6pp peak-to-final) when MUUFL tasks are introduced, a challenge shared by all tested methods. The method's effectiveness is tied to the backbone's spectral-spatial branch separation---a deliberate co-design choice validated by the ViT experiment (Section~\ref{sec:experiments}), which confirms that modality-aware architecture is essential for asymmetric adaptation.

\paragraph{Future work.}
Incorporating lightweight feature replay from stored statistics may further mitigate cross-domain forgetting. Extending the benchmark to more diverse sensor configurations (SAR, multispectral) and longer task sequences would test scalability.
